\section{Lecture 4}

\subsection*{Structure of Stationary Time Series}
\textbf{Def.} \(X_t\) is stationary if
\[(X_1, \dots, X_m) \overset{d}{=} (X_{1+\ell}, \dots X_{m + \ell})\]
We want to develop asymptotic theory for the time series \(X_t\).

\textbf{Def.} An ergodic stationary causal process can be written as
\[X_t = g(\epsilon_t, \epsilon_{t-1}, \dots)\]
where \(\epsilon_t\) is iid. This representation is due to Wiener (1956). He proved that every stationary, ergodic process can be written in this form.

\textbf{Example 1.} \(X_t = \sum_{i=0}^\infty a_i \epsilon_{t-i}\) and \(\sum_{i=0}^\infty a_i^2 , \infty\). This is called a \emph{linear process}.

\textbf{Example 2.} \(X_t = \epsilon_t \epsilon_{t-1}\).

\textbf{Example 3.} \(X_t = \sum_{0 \leq i \leq j} a_{ij} \epsilon_{t-i}\epsilon_{t-j}\) where \(\sum_{ij} |a_{ij}| < \infty\). This is called a \emph{Volterra process}. It arises naturally in the second-order Taylor expansion of the nonlinear filter. The generalization to higher order is called the \emph{Multiple Wiener-Ito Sum}.

\textbf{Example 4.} \(X_t = \rho |X_{t-1}| + \epsilon_t\). Then
\[g(\epsilon_t, \epsilon_{t-1}, \dots) = \epsilon_t + | \rho \epsilon_{t-1} + |\rho \epsilon_{t-2} + \cdots||\]
\textbf{Homework} Suppose \(\epsilon_t \sim N(0, 1)\). Show that the pdf of the above is skew-normal with density proportional to
\[c \phi(x)\Phi(\delta x), \quad \delta = \frac{\rho}{\sqrt{1 - \rho^2}}\]

\subsection*{Wiener\textquotesingle s Proof}
Let \(\mathcal F_t = \sigma(X_t, X_{t-1}, \dots).\)
\[[X_t \mid \mathcal F_{t-1}] \overset{d}{=} Q(U_t, \mathcal F_{t-1})\]
where \(Q\) is the conditional quantile function. Recall that a joint distribution \((X, Y)\) can be written as
\[(X, Y) \overset{d}{=} (Q(U \mid Y), Y),\]
where \(U\) and \(Y\) are independent. More broadly,
\[(X_t, \mathcal{F}_{t-1}) \overset{d}{=} (Q(U_t, \mathcal{F}_{t-1}), \mathcal{F}_{t-1})\]
where \(U_t \sim \text{Unif}[0, 1]\) iid independently of \(\mathcal{F}_{t-1}\). Next, we can write
\[\mathcal{F}_{t-1} \overset{d}{=} (X_{t-1}, \mathcal{F}_{t-2})\]
then plug in
\[\begin{gathered}X_t \overset{d}{=} Q(U_t, \mathcal{F}_{t-1}) \\
\overset{d}{=} Q(U_t, (X_{t-1}, \mathcal{F}_{t-2}))\end{gathered}\]
Then plug in \(X_{t-1} = Q(U_{t-1}, \mathcal{F}_{t-2})\) and repeat. Why does this argument not work? The resulting sequence of functions may not converge -\/- we need to impose some conditions.

\begin{center}\rule{0.5\linewidth}{0.5pt}\end{center}

\emph{Markov process}:

\(X_t = F(X_{t-1}, \epsilon_t)\) where \(\epsilon_t\) is iid. This is a special case where \(X_t \mid \mathcal{F}_{t-1}\) is equal to \(X_t \mid X_{t-1}\).

\(X_t = F(X_{t-1}, X_{t-2}, \epsilon_t)\) -\/- this is a \emph{second-order} Markov process. We can stack this to get a first-order Markvo process:
\[Y_t =
\begin{bmatrix}
X_t \\
X_{t-1}
\end{bmatrix}, \qquad Y_t = F(Y_{t-1}, \epsilon_t).\]
We can also call this a nonlinear AR process. We can also write this as
\[X_t = F_{\epsilon_t}(X_{t-1}), \qquad F_{\epsilon_t}: \R \rightarrow \R\]
as a system of iterated random functions. This is a simplification of the generic case. We can ask: under what conditions does this sequence of functions converge?

\textbf{Example 5.}
\[X_t = \epsilon_t \sqrt{a^2 + b^2 X_{t-1}^2}, \qquad a, b \in \R.\]
This is called *ARCH(1). The examples above have constant conditional variance (i.e. \(\text{Var}[X_t \mid \mathcal{F}_{t-1}]\)), but this model does not:
\[Var[X_t \mid X_{t-1} = x] = a^2 + b^2 x^2\]
This idea, due to Engle, explains volatility clustering well.

\textbf{Example 6.}
\[X_t = \epsilon_t \sqrt{a^2 + b^2 X_{t-1}^2 - c^2 X_{t-1}}\]
This is also a conditionally heteroskedastic process. However, this model is even more reasonablle than ARCH(1) to model volatility because it includes a directionality component \(X_{t-1}\) to capture volatility.

\textbf{Example 7.}
\[X_t =
\begin{cases}
    a_1 X_{t-1} + \epsilon_t, \quad X_{t-1} > 0 \\
    a_2 X_{t-2} + \epsilon_t, \quad X_{t-1} \leq 0.
\end{cases}\]
This is called a threshold AR process TAR(1). This was proposed by Howell Tong in the 1980s.

\textbf{Example 8.}
\[X_t = (\rho + b \epsilon_t) X_{t-1} + \epsilon_t\]
This is akin to an AR(1) process, with a random coefficient -\/- a \emph{random coefficient model}. This is use to model earthquakes.

\begin{center}\rule{0.5\linewidth}{0.5pt}\end{center}

\textbf{Question}: For Markov processes, can we find a distribution \(\pi\) such that \(X_t \Rightarrow \pi\)? This is slightly more general than stationarity.

Let\textquotesingle s revisit an old example. If
\[X_t = \rho X_{t-1} + \epsilon_t, \qquad F(x, \epsilon) = \rho x + \epsilon\]
Then we have convergence to stationarity if \(|\rho| < 1\). This motivates the general case: the map \(F(\cdot, \epsilon)\) should contract (on average?)

\textbf{Example 9.}

Consider \(X_t = F(X_{t-1})\). The Banach fixed-point theorem asserts that \(X_t \rightarrow x^*\) if \(F\) is a contraction mapping: \(\sup_x |F'(x)| < 1\).

However, we cannot have a fixed point in the random case because the map \(F(\cdot, \epsilon_t)\) changes with the random noise. We move from considering fixed points to considering fixed distributions: \(F(\pi) = \pi\).

Motivated by the proof of the Banach fixed-point theorem, we require
\[\mathbb{E} \left[\sup_{x \neq x'} \frac{F_{\epsilon_t}(x) - F_{\epsilon_t}(x')}{|x - x'|}\right] = \mathbb{E}[\text{Lip}(F_{\epsilon_t})] < 1.\]
If \(F\) is differentiable, we can require
\[\mathbb{E} \sup_x |F'_{\epsilon_t}(x)| < 1.\]
